\begin{figure*}[t!]
\centering
\includegraphics[width=0.91\linewidth]{./fig/model.pdf}
\caption{The overall framework of our ReSeq. Note that it only shows the single perspective of the model from $u_i$ to $v_j$, the other perspective from $v_j$ to $u_i$ is symmetrical.}
\label{fig:model}
\end{figure*}

\section{Methodology}
In this section, we first formulate our reciprocal sequential recommendation task, and secondly introduce the proposed ReSeq method in detail. The overall architecture is shown in Figure~\ref{fig:model}.


\subsection{Problem Formulation} \label{sec:pb_fmu}
In conventional recommender systems, the goal of recommendation algorithms is typically to predict the user's interest in items (movies, books, clothes, etc), which essentially refers to user-oriented item recommendation. 
Such a recommendation paradigm is solely based on the user perspective for item selection.  
As a comparison, Reciprocal Recommender System (RRS)~\cite{pizzato2010recon} targets a two-sided recommendation setting, in which either side can select the users or resources on the other side, 
\eg online dating recommendation between  females and males and job recommendation between job seekers and recruiters. 
Formally, we denote the two parties in reciprocal recommendation as $u_i \in \mathcal{U}$ (the first side)  and $v_j \in \mathcal{V}$ (the second side), respectively.
In this setting, we do not explicitly use the term of \emph{item}, since each side can play the roles of \emph{users} (those who actively select the interactions) and \emph{items} (those who are interacted with passively) as in traditional recommender systems. 
Instead, we call both sides \emph{users} and use different notations to discriminate the two sides. 
A successful recommendation in RRS indicates that: \textbf{(a)} the user $u_i$ demonstrates interest in the user $v_j$; \textbf{(b)} the features of $u_i$ align with the preferences and requirements of $v_j$. 
This process can be characterized as a dual-perspective active selection process~\cite{pizzato2010recon,xia2015reciprocal,neve2019latent,yang2022modeling}, the goal of RRS typically aims at learning a joint function $\hat{y}_{i,j}= \mathop{g}\limits^{\rightarrow}(u_i, v_j) +  \mathop{g}\limits^{\leftarrow}(u_i, v_j) $ that can  model this bilateral preference selection for the recommendation. 
In this paper, our approach is established in the aforementioned setting for RRS. 
This formulation characterizes a dual-perspective sequence matching task for the RRS, which is distinctive from previous settings~\cite{pizzato2010recon}.

To be more specific, at every time step, each user possesses a historical behavior sequence up to the current time $T$, denoted by $\mathcal{S}^T_{u_i} = (v^{t_1}_1, v^{t_2}_2,..., v^{t_a}_a)$ or $\mathcal{S}^T_{v_j} = (u^{t_1}_1, u^{t_2}_2,..., u^{t_b}_b)$ where $t_a < T$ and $t_b < T$ are the timestamps for interactions. In the sequence $\mathcal{S}^T_{u_i}$, $v^{t_{j'}}_{j'}$ represents that $u_i$ is matched with $v_{j'}$ at time $t_{j'}$. And vice versa, $u^{t_{i'}}_{i'}$ in $\mathcal{S}^T_{v_j}$ represents the historical matching relationship with $v_j$.
Our target is to construct a matching function $\hat{y}_{i,j}=\mathop{g}\limits^{\rightarrow}(u_i, {S}^T_{u_i}, v_j, {S}^T_{v_j}) + \mathop{g}\limits^{\leftarrow}(u_i, {S}^T_{u_i}, v_j, {S}^T_{v_j})$ that predicts the matching degree or score $\hat{y}_{i,j}$ between both users by dual-perspective modeling of their historical behavior sequences.

\subsection{Dual-Perspective Dynamic Sequential Behavior Modeling} 
\label{sec:bhv_model}

Existing RRS methods primarily utilize user profiles for prediction matching~\cite{pizzato2010recon,luo2020rrcn} or learning user static representations based on historical user interactions~\cite{xia2015reciprocal,neve2019latent}.
However, such methods tend to mainly focus on capturing long-term user interests, thereby ignoring the dynamic changes in user preferences over time.
Our approach models reciprocal recommendation as a distinctive sequence matching task, where the matching of both sides is conducted based on their dynamic behavior sequences.
In this section, we present our method for modeling user behavior sequences under dual perspectives.
To be more precise, we first illustrate the user dual-perspective representations in reciprocal recommendation, and then we elucidate the process of encoding both active and passive aspects of user behavior sequences.

\subsubsection{Bilateral User Representation Learning}
As mentioned in the previous section, the matching of users on both sides in reciprocal recommendation is a dual-perspective process.
This phenomenon results in the user playing a dual role in reciprocal recommendation, serving not only as the active selecting party but also as the passive accepting party.
Hence, methods in reciprocal recommendation frequently use dual-perspective user representations~\cite{neve2019latent,yang2022modeling}.
 
\paratitle{User representation learning under dual perspectives.}
For both users in the reciprocal recommendation scenario, there exist two distinct representations: one is when they act as active selectors by displaying their preferences over the other party, and the other is when they act as passive candidates by exhibiting their own features for being selected.
For simplicity, we use ``\emph{\underline{p}reference}'' as the marker for the user active representations and ``\emph{\underline{f}eature}'' as the marker for the user passive representations.
Particularly, for users $\mathcal{U}$ on one side, we maintain two separate embedding matrices $\bm{M}^{U}_{p} \in \mathbb{R}^{|\mathcal{U}| \times d}$ and $\bm{M}^{U}_{f} \in \mathbb{R}^{|\mathcal{U}| \times d}$, which represent user active and passive embeddings, respectively. Likewise, there are also two embedding matrices $\bm{M}^{V}_{p} \in \mathbb{R}^{|\mathcal{V}| \times d}$ and $\bm{M}^{V}_{f} \in \mathbb{R}^{|\mathcal{V}| \times d}$ for the users $\mathcal{V}$ on the other side.

\paratitle{Decomposed and shared bilateral user embedding.}
As discussed previously, the matching of both parties in reciprocal recommendation is a dual-perspective alignment of their active interests and passive attractions.
This implies that the representation space of bilateral users should be properly aligned.
Formally, the above embedding matrices are aligned between $\bm{M}^{U}_{p}$ and $\bm{M}^{V}_{f}$, as well as between $\bm{M}^{V}_{p}$ and $\bm{M}^{U}_{f}$.
To reflect this dual-perspective alignment relationship, we perform a decomposition of the user embedding matrix, and part of the decomposed matrix is shared between the dual perspectives. Take a perspective as an example:
\begin{align}
    \bm{M}^{U}_{p} = \bm{A} \cdot \bm{C}, \ \ \ \ \
    \bm{M}^{V}_{f} = \bm{B} \cdot \bm{C},
    \label{eq:decomp}
\end{align}
where $\bm{A} \in \mathbb{R}^{|\mathcal{U}| \times {d'} }$, $\bm{B} \in \mathbb{R}^{|\mathcal{V}| \times {d'}}$, and $\bm{C} \in \mathbb{R}^{{d'} \times d}$ is shared between the user active representations on one side and the user passive representations on the other side. Another perspective also applies the identical approach of matrix decomposition and sharing.


\subsubsection{Dual-Perspective Behavior Sequence Encoding}
\label{sec:seq_encode}
After obtaining the decomposed and shared dual-perspective user embedding, we next encode the user behavior sequence by proposing two different methods. The purpose is to actively and passively model behavior sequences according to the different dynamic characteristics of the two roles for the user.

\paratitle{Unidirectional active dynamic encoding.}
When a user assumes the role of the active party, each historical matching user in the behavior sequence acts as the passive party.
Intuitively, the features and characteristics of passive candidates in recommendations are indicative of user interests and preferences.
Therefore, we encode the active dynamic representations of the user based on the passive representations of historical matching users in the behavior sequence.

Given a behavior sequence $\mathcal{S}^T_{u_i} = (v^{t_1}_1, v^{t_2}_2,..., v^{t_a}_a)$, we first pad the sequence to a fixed length $n$, and then add a special token ($[CLS]$) at the beginning of the sequence to aggregate the sequence representations. 
Next, we apply the lookup operation from $\bm{M}^{V}_{f}$ along with a learnable position embedding $P \in \mathbb{R}^{(n+1) \times d}$ to form the input sequence passive embeddings $\bm{E}^{V}_{f} \in \mathbb{R}^{(n+1) \times d}$. 
Subsequently, we exploit a widely-used Transformer network~\cite{vaswani2017attention} for sequence modeling. Concretely, a typical Transformer architecture consists of multiple layers of masked multi-head self-attention (denoted by Masked-MHAttn(·)), and each layer of self-attention followed by a point-wise feed-forward network (denoted by FFN(·)). The above sequence embeddings $\bm{E}^{V}_{f}$ combined with position encoding are used as the initial input of the Transformer. Then, the calculations can be formally written as:
\begin{align}
    \bm{H}^{l} & =\operatorname{FFN}\left(\operatorname{Masked-MHAttn}\left(\bm{H}^{l-1}, \bm{W}_{mask}\right)\right), \ l \in\{1, \ldots, L\},
\end{align}
where $\bm{H}^{l}$ is the hidden states in the $l$-th layer, and $\bm{H}^{0} = \bm{E}^{V}_{f}$. $\bm{W}_{mask} \in \mathbb{R}^{(n+1) \times (n+1)}$ is the attention mask weight matrix used to regulate the attention object for each part within the sequence. 

In our case, to encode the time-sensitive active dynamic behaviors, we use a unidirectional attention mask weight matrix which allows the model to concentrate more on recent user behaviors. 
As shown in Figure~\ref{fig:model}, compared with the diagonal matrix unidirectional mask, our mask matrix separately processes the corresponding row and column of the $[CLS]$ token to ensure both the order of information transmission and the global aggregation of $[CLS]$ representation. 
Note that the mask matrix used in our approach is equivalent to adding the $[CLS]$ token at the end of the sequence. However, compared to individually inserting $[CLS]$ token to behavior sequences of different lengths, our approach is more convenient for input construction and model implementation.

Ultimately, the outputs from the final layer of the $L$-layer Transformer network serve as the active dynamic representations of the user. Additionally, we divide $\bm{H}^{L}$ into two scales:
\begin{align}
    \bm{p}_{u_i} = \bm{H}^{L}_{0}, \ \ \ \ \
    \widetilde{\bm{p}}_{u_i} = \bm{H}^{L}_{1:n+1},
\end{align}
where $\bm{H}^{L}_{0}$ is the corresponding output of $[CLS]$ token, $\bm{p}_{u_i} \in \mathbb{R}^{d}$ indicates the user's macro or global active dynamic representation. And $\bm{H}^{L}_{1:n+1}$ is the other part of the outputs, $\widetilde{\bm{p}}_{u_i} \in \mathbb{R}^{n \times d}$ represents the user's micro or fine-grained active dynamic representations.

\paratitle{Bidirectional passive dynamic encoding.}
Similar to the unidirectional active dynamic encoding described above, we also employ a Transformer network to model the passive perspective of user behavior sequences.
The difference is that, in this context, each historical matching user in the behavior sequence assumes the role of the active party.
Empirically, a successful matching means that the features and characteristics of the current user meet the preferences and requirements of the active selector.
We utilize the lookup operation from $\bm{M}^{V}_{p}$ instead of $\bm{M}^{V}_{f}$ to form the input sequence active embeddings $\bm{E}^{V}_{p} \in \mathbb{R}^{(n+1) \times d}$.

Furthermore, in contrast to the time-sensitivity of user active dynamic behaviors, when a user acts as a passive individual, his features and characteristics tend to remain relatively stable over time.
Thus, as shown on the right side of Figure~\ref{fig:model}, we encode the passive dynamic representations by a bidirectional attention mask, which enables the model to capture essential user traits across global sequence representations. 
In the same way, the outputs of the last layer are also divided into two scales:
\begin{align}
    \bm{f}_{u_i} = \bm{H}^{L}_{0}, \ \ \ \ \
    \widetilde{\bm{f}}_{u_i} = \bm{H}^{L}_{1:n+1},
\end{align}
where $\bm{f}_{u_i} \in \mathbb{R}^{d}$ is the user's macro passive dynamic representation, and $\widetilde{\bm{f}}_{u_i} \in \mathbb{R}^{n \times d}$ is the user's micro passive dynamic representations.

For user $v_{j}$ on the other side, we also encode his active and passive dynamic representations based on the symmetrical model structure, which are denoted by $\bm{p}_{v_j}$, $\widetilde{\bm{p}}_{v_j}$ and $\bm{f}_{v_j}$, $\widetilde{\bm{f}}_{v_j}$.

\subsection{Multi-Scale Sequence Matching} \label{sec:match}
The above process allows us to obtain multi-scale active and passive dynamic representations of users by encoding their behavior sequences. In the following section, we will delve into the dual-perspective multi-scale matching of both users in reciprocal recommendation.
\subsubsection{Macro-Level Matching}
As mentioned in the previous sections, the matching of both parties in reciprocal recommendation is a dual-perspective active selection process~\cite{pizzato2010recon,xia2015reciprocal,neve2019latent,yang2022modeling}. Hence, we apply two scores for active-passive matching to model this dual-perspective propensity. It is worth noting that in Figure~\ref{fig:model}, only one perspective matching process is depicted, while the other perspective is symmetrical. The matching prediction method at the macro level is as follows:
\begin{align}
\begin{split}
    {\hat{y}}_{i,j} & = {\hat{y}}_{i \rightarrow j} + {\hat{y}}_{j \rightarrow i}\\
    &= (\bm{p}_{u_i})^{T} \cdot \bm{f}_{v_j} + (\bm{p}_{v_j})^{T} \cdot \bm{f}_{u_i},
\end{split}
\label{eq:ma_score}
\end{align}
where $\cdot$ is the dot product operation. The first part of the formula represents the propensity from $u_i$ to $v_j$, and the second part represents the propensity from $v_j$ to $u_i$.
\subsubsection{Time-Sensitive Micro-Level Matching}
Compared with the conventional recommender systems that rely on a fixed number of static items as candidate-matching targets, our problem formulation involves a distinctive sequence matching task, which employs dynamic user behavior sequences as candidate-matching targets. This results in a broader and more dispersed candidate matching space. 
Consequently, to alleviate the sparsity of the candidate matching space, we strive to attain more discriminative and precise matching predictions by incorporating more detailed and fine-grained micro-level interactions between users on both sides. 
Formally, the overall micro-level matching method is as follows:
\begin{align}
    {\hat{z}}_{i,j} = \operatorname{TiSensiMatch} \left(\widetilde{\bm{p}}_{u_i}, \ \widetilde{\bm{f}}_{v_j}\right) + \operatorname{TiSensiMatch} \left(\widetilde{\bm{p}}_{v_j}, \  \widetilde{\bm{f}}_{u_i}\right).
    \label{eq:mi_score}
\end{align}

Next, as shown in the ``Time-Sensitive Micro-Level Matching'' section of Figure~\ref{fig:model}, we take $\operatorname{TiSensiMatch} \left(\widetilde{\bm{p}}_{u_i}, \ \widetilde{\bm{f}}_{v_j}\right)$ as an example to describe the implementation of time-sensitive matching in detail. 
Firstly, we compute the fine-grained matching matrix between the micro active dynamic representations $\widetilde{\bm{p}}_{u_i}$ and the passive dynamic representations $\widetilde{\bm{f}}_{v_j}$:
\begin{align}
    \bm{G} = \widetilde{\bm{p}}_{u_i} \cdot {(\widetilde{\bm{f}}_{v_j})}^{T},
\end{align}
where $\bm{G} \in \mathbb{R}^{n \times n}$ stands for the matching degrees between the active and passive representations at different time steps. 
Secondly, we aggregate the matrix $\bm{G}$ through two dimensions of co-attention. For the passive dimension, we compute attention weights based on softmax:
\begin{align}
    \bm{\gamma} = \operatorname{softmax}(\widetilde{\bm{f}}_{v_j} \cdot \bm{e}^{p}_{u_i}),
\end{align}
where $\bm{e}^{p}_{u_i} \in \mathbb{R}^{d}$ is the original active embedding of the user $u_i$ from $\bm{M}^{U}_{p}$. For the active dimension, we propose a relative time-dependent attention weight to model the time-sensitivity of user active dynamic behaviors:
\begin{align}
    \bm{\delta} = \operatorname{softmax}(\widetilde{\bm{p}}_{u_i} \cdot \bm{e}^{f}_{v_j} + \bm{\alpha}^{\tau}),
\end{align}
where $\bm{e}^{f}_{v_j} \in \mathbb{R}^{d}$ is the original passive embedding of the user $v_j$ from $\bm{M}^{V}_{f}$. 
And $\bm{\alpha}^{\tau}$ is from a learnable relative time weight $\bm{\alpha} \in \mathbb{R}^{n}$ starting from the end of the sequence. 
For a behavior sequence with an actual length of $s$ (no padding), there is $\bm{\alpha}^{\tau}_{s-k} = \bm{\alpha}_{k-1}$, where $k\in\{1,2,\ldots,s\}$. 
Finally, we aggregate the fine-grained matching matrix $\bm{G}$ in the following way:
\begin{align}
\begin{split}
     {\hat{z}}_{i \rightarrow j} &= \operatorname{TiSensiMatch} \left(\widetilde{\bm{p}}_{u_i}, \ \widetilde{\bm{f}}_{v_j}\right) \\
    & =  {\bm{\delta}}^{T} \cdot \bm{G} \cdot \bm{\gamma}.
\end{split}
\label{eq:agg}
\end{align}

In addition, for the perspective from $v_j$ to $u_i$, we also get the micro propensity score ${\hat{z}}_{j \rightarrow i}$ in the same way. The full micro-level matching score ${\hat{z}}_{i,j}$ is the sum of the dual-perspective propensities.

\subsection{Improving Matching Efficiency via Self-Distillation}
The rapid development of the internet results in an increased number of users on online platforms. However, this also leads to a service efficiency problem that needs to be solved urgently. 
Specific to the multi-scale sequence matching method we proposed, macro-level matching is based on two simple and efficient dot product operations, while micro-level matching provides more fine-grained bilateral interactions. 
Nevertheless, the complexity of the latter also presents efficiency issues that need to be addressed.

Inspired by knowledge distillation techniques~\cite{hinton2015distilling} that have been successfully applied in various fields~\cite{zhang2019your,sanh2019distilbert,hofstatter2020improving,tian2023directed}, we also aspire to improve the recommendation efficiency in practical applications through self-distillation from micro-level matching to macro-level matching.
More exactly, as shown in the ``Micro-to-Macro Self-Distillation'' section of Figure~\ref{fig:model}, we use the self-distillation technique to transfer knowledge from finer-grained micro-level matching to more efficient and simple macro-level matching during training. 
Furthermore, only macro-level matching is used in the process of verification, testing and practical application to achieve efficient and precise recommendations.

\paratitle{Micro-to-macro self-distillation.}
As a ranking task, reciprocal recommendation differs from common classification problems in that it pays more attention to the relative score gap between positive and negative instances rather than the absolute category distribution. Therefore, we adopt Margin-MSE~\cite{hofstatter2020improving} as our distillation loss to facilitate the alignment of positive and negative score margins between micro-level and macro-level matching. Specifically, the calculation for self-distillation loss is:
\begin{align}
\begin{split}
\mathcal{L}_{sd}
    & = \operatorname{Margin-MSE} \left({\hat{z}}^{+}, \ {\hat{z}}^{-}, \ {\hat{y}}^{+}, \ {\hat{y}}^{-}\right), \\
    & = \frac{1}{|B|} \sum^{|B|}_{i=1} \left[(\hat{z}^{+}_{i} - \hat{z}^{-}_{i}) - (\hat{y}^{+}_i -\hat{y}^{-}_i )\right]^{2},
\end{split}
\end{align}
where $|B|$ represents the batch size of training data. Besides, ${\hat{z}}^{+}$, ${\hat{z}}^{-}$ and ${\hat{y}}^{+}$, ${\hat{y}}^{-}$ are the positive and negative score sets of micro level and macro level, respectively. 

\paratitle{Overall optimization.}
In our approach, the final matching prediction is computed by four elements from two users (Eqn.~\eqref{eq:ma_score} and ~\eqref{eq:mi_score}), which are indispensable in the actual scenario for successful matching. Thus, we can get four negative instance scores by replacing the four elements with negative instances, respectively. Formally, for each positive instance $(u_i, {S}^T_{u_i}, v_j, {S}^T_{v_j})$, we separately sample bilateral negative users $(u_{i'}, {S}^T_{u_{i'}})$ and $(v_{j'}, {S}^T_{v_{j'}})$. Note that the historical behavior sequences of negative users are also truncated to the current time $T$. After inputting negative instances into the model, four sets of negative instances at the macro level and micro level can be obtained:
\begin{align}
    \bm{D} = \{(\bm{p}_{u_{i'}}, \widetilde{\bm{p}}_{u_{i'}}),(\bm{f}_{u_{i'}}, \widetilde{\bm{f}}_{u_{i'}}),(\bm{p}_{v_{j'}}, \widetilde{\bm{p}}_{v_{j'}}),(\bm{f}_{v_{j'}}, \widetilde{\bm{f}}_{v_{j'}})\}.
\end{align}

\textcolor{black}{By replacing the corresponding elements in Eqn.~\eqref{eq:ma_score} and ~\eqref{eq:mi_score} with the aforementioned negative instances, we can compute multiple negative instance scores for both scales. Without loss of generality, we denote them as two negative instance score sets, $\bm{Y}^{-}_{i,j}$ and $\bm{Z}^{-}_{i,j}$, each containing four types of negative instance scores.} 
Then, we use the BPR loss~\cite{rendle2009bpr} to optimize the ranking between positive and negative instances:
\begin{align}
    \mathcal{L}_{ma}
    & = - \frac{1}{|B|} \sum_{(u_i,v_j) \in B} \ \sum_{\hat{y}^{-}_{i,j} \in \bm{Y}^{-}_{i,j}} \log \left( \sigma ( \hat{y}^{+}_{i,j}-\hat{y}^{-}_{i,j})\right), \\
    \mathcal{L}_{mi}
    & = - \frac{1}{|B|} \sum_{(u_i,v_j) \in B} \ \sum_{\hat{z}^{-}_{i,j} \in \bm{Z}^{-}_{i,j}} \log \left( \sigma ( \hat{z}^{+}_{i,j}- \hat{z}^{-}_{i,j} )\right),
\end{align}
where $B$ represents a batch of training data. $\hat{y}^{+}_{i,j}$ refers to the macro-level matching score of user pairs, and $\hat{y}^{-}_{i,j}$ denotes its negative instance score. $\hat{z}^{+}_{i,j}$ and $\hat{z}^{-}_{i,j}$ are corresponding to micro-level matching. In addition, based on multiple negative instance scores, the self-distillation loss can also be extended to:
\begin{align}
    \mathcal{L}_{sd}
     = \frac{1}{|B|} \sum_{(u_i,v_j) \in B}  \sum_{\scriptscriptstyle{\hat{y}^{-}_{i,j} \in \bm{Y}^{-}_{i,j}, \hat{z}^{-}_{i,j} \in \bm{Z}^{-}_{i,j}}} \scriptstyle{\left[(\hat{z}^{+}_{i,j} - \hat{z}^{-}_{i,j}) - (\hat{y}^{+}_{i,j} -\hat{y}^{-}_{i,j})\right]^{2}}.
     \label{eq:lsd}
\end{align}


Combining the above three losses, our overall optimization loss function can be denoted as follows: 
\begin{align}
    \mathcal{L} = \mathcal{L}_{ma} + \lambda \mathcal{L}_{mi} + \mu \mathcal{L}_{sd},
\end{align}
where $\lambda$ and $\mu$ are hyperparameters. $\lambda$ is used to balance the training of macro and micro, while $\mu$ is for the trade-off between ranking optimization and self-distillation knowledge transfer.

Note that during validation and testing, we only deploy macro-level matching module, which makes the model much more efficient in practice.

\subsection{Discussion}
\subsubsection{Time Complexity Analysis}
In a real-world deployment, the model inference associated only with the entities themselves can be completed in advance.
In reference to our approach, the encoding of a user's behavior sequence is independent of other individuals, and the outcomes can be pre-calculated and stored. 
Therefore, we mainly discuss the time complexity of the matching prediction part.
Matching at the macro level requires only the time complexity of $\mathcal{O}(d)$, where $d$ denotes the model output dimension. It is the same as many classical collaborative filtering-based methods~\cite{rendle2009bpr,neve2019latent} and sequential recommendation methods~\cite{kang2018self,sun2019bert4rec}.
For micro-level matching, it first takes $\mathcal{O}(n^2d)$ time to complete the calculation of the fine-grained matching matrix, where $n$ denotes the sequence length. Then, the time complexity for the co-attention of the two dimensions is $\mathcal{O}(nd)$. The final aggregation operation still takes time of $\mathcal{O}(n^2)$. Overall, the time complexity of prediction at the micro level is $\mathcal{O}(n^2d)$.
Given that online platforms cater to tens of thousands of users, such time complexity is unacceptable. As a result, we suggest utilizing self-distillation from the micro level to the macro level to optimize the model's efficiency for practical application.
\subsubsection{Comparison with Existing Recommendation Methods}
Here, we briefly compare with related recommendation methods to highlight the innovation and novelty of our approach.

\paratitle{Sequential recommendation methods.}
Existing sequential recommendation algorithms such as SASRec~\cite{kang2018self} and BERT4Rec~\cite{sun2019bert4rec} only encode user dynamic representations from one perspective, which lacks the overall modeling of dual perspectives in the reciprocal recommendation scenario. In contrast, the methodology we proposed involves the dual-perspective joint modeling of users. And we achieve more precise recommendations by performing fine-grained matching on more details of both recommending parties.

\paratitle{Reciprocal recommendation methods.}
In the current reciprocal recommendation studies (regardless of content-based or collaborative filtering methods), the modeling of the user ultimately results in a static representation vector. 
Some methods~\cite{pizzato2010recon,yildirim2021bideepfm} model user preferences based on their fixed profiles, while others~\cite{rendle2009bpr,neve2019latent} learn user's static implicit representations through matrix factorization and latent factor models. 
Even methods that are also based on user historical behaviors~\cite{fu2021beyond} only aggregate all user histories in the training set, while ignoring the temporal dynamics of behavior sequences. 
As a comparison, our method takes into full consideration the dynamics on both user active and passive perspectives when constructing model inputs (Section~\ref{sec:pb_fmu}), encoding user behaviors (Section~\ref{sec:bhv_model}), and predicting bilateral user matching (Section~\ref{sec:match}).











